% =============================================================================
% Chapter 1: General Project Presentation  (6–8 pages target)
% =============================================================================
\chapter{General Project Presentation}
\label{chap:presentation}

\section*{Introduction}

This chapter introduces the overall context and motivation behind the \eduspace{} project.
It begins by describing the problem domain---the shortcomings of current classroom management platforms---and states the specific objectives that the project seeks to achieve.
It then explains how the team organized its work, the development methodology adopted, and concludes with a comparative study of existing solutions to justify the design decisions made.

% ─────────────────────────────────────────────────────────────────────────────
\section{Context and Problem Statement}
\label{sec:context}

The digitization of higher education has progressed substantially over the past decade.
Learning Management Systems (LMS) and classroom management platforms have become indispensable tools for organizing course materials, distributing assignments, tracking student progress, and facilitating communication between instructors and learners.
Platforms such as Google Classroom~\cite{googleclassroom2024}, Moodle~\cite{moodle2024}, and Microsoft Teams for Education~\cite{msteams2024edu} are now deployed at scale across universities worldwide.

Despite their widespread adoption, these tools exhibit recurring limitations that hinder the learning experience:

\begin{enumerate}[label=(\roman*)]
  \item \textbf{Poor content organization.}
        In platforms like Google Classroom, all course materials---lectures, exercises, lab instructions, summaries---are displayed on a single, flat feed sorted by date.
        As the volume of content grows, students find it increasingly difficult to locate specific resources.
        There is no categorization by type (lecture vs.\ lab vs.\ exercise) or by chapter, and no dedicated search functionality within a classroom.

  \item \textbf{Limited communication channels.}
        Most platforms lack a built-in, dedicated space for student-to-student interaction.
        While some offer comment threads on posts, they do not provide real-time group chat.
        As a result, students typically create informal groups on external social media applications (e.g., Facebook Messenger, WhatsApp, Telegram), which introduces distractions, fragments the learning workflow, and excludes students who do not use a particular network.

  \item \textbf{Rigid classroom models.}
        Existing platforms assume a single classroom paradigm---typically a formal teacher--student hierarchy.
        They do not natively support informal or peer-led study groups where grading is irrelevant and all participants share equal standing.

  \item \textbf{Lack of advanced search.}
        Full-text search across post titles, descriptions, and attachment file names within a specific classroom is absent in most platforms, forcing users to scroll through chronological feeds manually.
\end{enumerate}

These shortcomings motivated the design of \eduspace{}: a classroom management platform that specifically targets content organization, in-platform real-time communication, flexible classroom types, and full-text search.

% ─────────────────────────────────────────────────────────────────────────────
\section{Specific Project Objectives}
\label{sec:objectives}

The primary goal of this project is to design, develop, and deploy a fully functional web-based classroom management platform that improves upon existing solutions in terms of organization, communication, and flexibility.
The specific objectives are:

\begin{enumerate}[label=\textbf{O\arabic*.}]
  \item \textbf{Structured content management.}
        Provide a system where every post belongs to a chapter and has an explicit type (Study Material, Quiz, Question, Assignment, or Announcement).
        Study materials are further categorized into sub-types: Cours, TD (tutorial exercises), TP (lab work), and Résumé (summary).
        Users can filter and sort the classroom stream by type, sub-type, chapter, and date.

  \item \textbf{Dual classroom model.}
        Support two distinct classroom modes---\emph{Teaching} (formal, with grading) and \emph{Friendly} (informal, without grading)---while reusing the same underlying data model and role system (Admin / Member).

  \item \textbf{Integrated real-time group chat.}
        Embed a WebSocket-based group chat within each classroom, with typing indicators, online presence tracking, file sharing, and paginated message history.
        The chat can be enabled or disabled by classroom administrators.
\end{enumerate}

% ─────────────────────────────────────────────────────────────────────────────
\section{Work Organization and Task Distribution}
\label{sec:organization}

The project was carried out by a team of two developers.
Rather than assigning rigid, technology-based roles (e.g., one person responsible for the entire frontend and another for the entire backend), the team adopted a \textbf{feature-based collaboration model}.
Each team member assumed full end-to-end ownership of a given feature---designing the data model, implementing the backend logic and API endpoints, and building the corresponding frontend interface.

This approach ensured that both developers maintained a holistic understanding of the system, from database schema to user interface.
Cross-cutting concerns such as UML diagram design, shared utility development, and architectural decisions were handled collaboratively.

Source code was managed using \textbf{Git} with a remote repository hosted on \textbf{GitHub}.
Feature branches were used to isolate work on individual services or features, and changes were integrated into the main branch through pull requests, enabling peer review before merging.
This workflow provided traceability, facilitated parallel development, and minimized integration conflicts.

% ─────────────────────────────────────────────────────────────────────────────
\section{Adopted Development Approach}
\label{sec:methodology}

Given the two-person team size and the microservices nature of the project, the team adopted an \textbf{iterative and incremental} development approach rather than a prescriptive framework such as Scrum.

Work was organized by \textbf{service and feature}: each microservice was developed end-to-end in a single iteration---backend logic, API endpoints, database schema, and the corresponding frontend views.
Once a service was functional and tested in isolation, the team proceeded to the next.
Cross-service integration (event wiring, inter-service REST calls) was addressed in dedicated integration iterations.

Each service followed the same iterative cycle: database schema design, backend logic and API endpoints, frontend UI, service-level testing, and cross-service integration---before proceeding to the next service.

% ─────────────────────────────────────────────────────────────────────────────
\section{State of the Art}
\label{sec:state-of-art}

% Note: Section 1.5.1 "Collecte de l'information" is facultatif — omitted per instructions.

\subsection{Existing Solutions}
\label{sec:existing-solutions}

Before defining the feature set of \eduspace{}, an analysis of the most prominent classroom management platforms was conducted.
The three solutions examined are Google Classroom, Moodle, and Microsoft Teams for Education.

\medskip
\tabref{tab:comparison} provides a comparative summary of these platforms against the key criteria addressed by \eduspace{}.

\begin{table}[H]
  \centering
  \caption{Comparison of existing classroom management platforms.}
  \label{tab:comparison}
  \small
  \begin{tabularx}{\textwidth}{@{} l *{4}{>{\centering\arraybackslash}X} @{}}
    \toprule
    \textbf{Criterion} & \textbf{Google Classroom} & \textbf{Moodle} & \textbf{MS Teams Edu} & \textbf{EduSpace} \\
    \midrule
    Content categorization by type     & No  & Partial & No  & Yes \\
    Chapter-based organization         & No  & Yes     & No  & Yes \\
    Full-text search within classroom  & No  & Partial & No  & Yes \\
    Integrated real-time group chat    & No  & No      & Yes & Yes \\
    Dual classroom model (formal/informal) & No & No   & No  & Yes \\
    Assignment \& grading workflow     & Yes & Yes     & Yes & Yes \\
    File attachments on all post types & Yes & Yes     & Yes & Yes \\
    Real-time notifications            & Partial & No  & Yes & Yes \\
    Open-source                        & No  & Yes     & No  & Yes \\
    Lightweight deployment (Docker)    & N/A & Complex & N/A & Yes \\
    Modern, responsive UI              & Yes & No      & Yes & Yes \\
    \bottomrule
  \end{tabularx}
\end{table}

% ─────────────────────────────────────────────────────────────────────────────
\subsection{Critiques and Recommendations}
\label{sec:critiques}

The comparative analysis reveals several gaps that no single existing platform fully addresses:

\begin{enumerate}[label=(\roman*)]
  \item \textbf{Content discoverability.}
        Google Classroom and Microsoft Teams for Education both lack structured content categorization.
        While Moodle offers course sections, its interface complexity offsets this advantage.
        \emph{Recommendation:} Combine type-based and chapter-based organization with full-text search powered by a dedicated search engine.

  \item \textbf{In-platform communication.}
        Only Microsoft Teams provides real-time chat, but it is embedded within a large, general-purpose collaboration suite.
        Google Classroom and Moodle offer no real-time messaging.
        \emph{Recommendation:} Integrate a lightweight, WebSocket-based group chat directly within each classroom, with the option for administrators to enable or disable it.

  \item \textbf{Classroom flexibility.}
        All three platforms assume a single classroom paradigm.
        \emph{Recommendation:} Support both a formal (teaching) mode with full grading capabilities and an informal (friendly) mode for peer-led study groups, reusing the same underlying architecture.
  \end{enumerate}

These recommendations directly informed the design decisions and feature set of \eduspace{}, as detailed in the subsequent chapters.

% ─────────────────────────────────────────────────────────────────────────────
\section*{Conclusion}

This chapter established the context and motivation for \eduspace{} by identifying the organizational and communicational limitations of existing classroom management platforms.


The next chapter details the requirements analysis, including actor identification, functional and non-functional requirements, and UML modeling of the system's behavior.
